\chapterhead{5}{Separating Everything}{Putting the tools in the right order}\label{ch5}

Now we can use all the tools together. Here is a mixture of three parts: steel paper clips, sand and salt. No single tool can take it apart, but the tools one after another can.

\begin{center}
\begin{tikzpicture}[font=\sffamily\scriptsize,node distance=0.4cm]
  \node[bigcir,minimum size=1.9cm] (m) at (0,0) {mixture 80};
  \node[cir] (c) at (3.2,1.4) {10}; \node[note,anchor=west] at (3.7,1.4) {clips (magnet)};
  \node[bigcir,minimum size=1.7cm] (r) at (3.4,-0.4) {70};
  \node[cir] (s) at (7.2,0.6) {49}; \node[note,anchor=west] at (7.7,0.6) {sand (filter)};
  \node[cir] (t) at (7.2,-1.6) {19}; \node[note,anchor=west] at (7.7,-1.6) {salt (evaporate)};
  \draw[popA] (m) -- node[above left,note] {magnet} (c);
  \draw[popA] (m) -- (r);
  \draw[popA] (r) -- node[above,note] {water, stir, filter} (s);
  \draw[popA] (r) -- node[below,note] {evaporate} (t);
\end{tikzpicture}
\end{center}

\begin{tryit}{Sand, salt and clips}
\textbf{You need:} a kitchen scale, a magnet, a filter, a saucer, three cups, water, 10~g of paper clips, 50~g of sand and 20~g of salt.
\begin{enumerate}[leftmargin=1.6em]
\item Mix the clips, sand and salt. Weigh the mixture.
\item Pull the clips out with the magnet. Weigh the clips and what is left.
\item Stir the rest into 100~g of water so that the salt dissolves. The sand will not dissolve.
\item Filter it. The sand stays on the filter. Dry the sand on a plate and weigh it.
\item Pour the salty water that came through onto a weighed saucer and let it evaporate. Weigh the salt crust.
\end{enumerate}
\par\medskip
\renewcommand{\arraystretch}{1.6}
\begin{tabular}{|l|c|c|c|c|c|}
\hline
 & \textbf{Whole} & \textbf{Clips} & \textbf{Sand} & \textbf{Salt} & \textbf{Parts added}\\ \hline
\textbf{Weight (g)} & & & & & \\ \hline
\end{tabular}
\end{tryit}
\adult{Do not taste anything. Keep small magnets and paper clips away from young children and pets. Wash hands after handling sand and salt water.}

\begin{worked}{Practice numbers, not from a real test}
Whole: $10+50+20=80$~g. After the magnet: clips 10~g, left 70~g.\\
Sand, dried: 49~g. Salt crust: 19~g.\\
Parts added: $10+49+19=78$~g. The gap is 2~g. We weighed three separate parts, so the scale's 1~g error could add up to 3~g, and a 2~g gap fits.
\end{worked}

\begin{multix}[hmm]\multixbot[hmm]\tcblower
I will filter the dry mixture first, because the filter will keep the sand and the salt together. Then I will separate them. That is quicker, no? Did I miss something?
\end{multix}

\trust{We weighed the whole mixture and each part, and we did the steps in an order where each tool had a difference to use.}
{The parts added to 78~g against 80~g, a gap smaller than the combined scale error.}
{That a bigger gap would be harmless, or that we can tell which step lost what from only the final weights. The next chapter shows how to find the step.}

\begin{keepbox}
We separated three parts. \textbf{We kept} \wline[3.4cm] \textbf{and we set aside} \wline[3.4cm].
\end{keepbox}

\clearpage

\begin{puzzle}{Order the steps}
Number these from 1 to 6.\\
\wline[1cm] Add water and stir. \quad \wline[1cm] Weigh the whole mixture. \quad \wline[1cm] Evaporate the salt water.\\
\wline[1cm] Pull out the clips with a magnet. \quad \wline[1cm] Filter the mixture. \quad \wline[1cm] Weigh the dry sand.
\end{puzzle}

\begin{puzzle}{After the magnet}
The mixture weighs 80~g. The magnet picks up 10~g of clips. How much is left? \wline[1.6cm]~g
\end{puzzle}

\begin{puzzle}{Add the parts}
Clips 10~g, sand 49~g, salt 19~g. Total \wline[1.6cm]~g. Whole mixture 80~g. Gap \wline[1.6cm]~g.
\end{puzzle}

\begin{puzzle}{Why dissolve first?}
Why must the salt dissolve before filtering? \wline[6cm]
\end{puzzle}

\begin{puzzle}{Which tool?}
For each part, write the tool: magnet, filter or evaporation.\\
Clips: \wline[3cm] \quad Sand: \wline[3cm] \quad Salt: \wline[3cm]
\end{puzzle}

\clearpage

\begin{puzzle}{Where did the sand lose a gram?}
You started with 50~g of sand and ended with 49~g. Give two places the missing gram could be. \wline[6cm]
\end{puzzle}

\begin{puzzle}{Where did the salt lose a gram?}
Starting with 20~g of salt, the crust weighed 19~g. Give two places the missing gram could be. \wline[6cm]
\end{puzzle}

\begin{puzzle}{A bigger gap}
Another group's parts added to 71~g from a mixture of 80~g. Is a 9~g gap within error? \wline[2cm] What would you do? \wline[4cm]
\end{puzzle}

\begin{puzzle}{Rice, salt and clips}
Plan how to separate rice, salt and clips. List the steps in order.\\
\wline[9cm]\\[8pt]\wline[9cm]
\end{puzzle}

\begin{puzzle}{Spot the unfair test}
Group A dried their sand for one hour and Group B for a whole day. Group B's sand weighed 3~g less. Why? \wline[6cm]
\end{puzzle}
\clearpage

\thispagestyle{fancy}
\subhead{Plan your own separation}
Pick a mixture of three parts, such as rice, salt and paper clips, or lentils, sand and salt. Plan the steps before you start, then record what happened.\par\medskip
\textbf{My mixture and its parts:} \wline[7cm]\\[10pt]
\textbf{The difference each tool will use:}\\[4pt]
\wline[9cm]\\[8pt]\wline[9cm]\\[8pt]
\renewcommand{\arraystretch}{2.6}
\begin{tabularx}{\linewidth}{|>{\sffamily\small\bfseries}p{1.6cm}|X|X|X|}
\hline
Step & What I will do & Expected (g) & Weighed (g) \\ \hline
1 & & & \\ \hline
2 & & & \\ \hline
3 & & & \\ \hline
4 & & & \\ \hline
5 & & & \\ \hline
\end{tabularx}

\vspace{8pt}
Total gap: \wline[2cm]~g \quad Where do I think it went? \wline[5cm]
\clearpage
